\شروع‌مسئله[40] \متن‌سیاه{پیشنهاد الگوریتم یا داده‌ساختار مناسب}. در هر قسمت از این سؤال از شما خواسته می‌شود یک الگوریتم یا داده‌ساختار مناسب پیشنهاد دهید. این داده‌ساختار یا الگوریتم می‌تواند از مطالب گفته شده در درس، ترکیبی از آن‌ها و یا خلاقیتی باشد. علاوه بر \زیرخط{کلیّت داده‌ساختار و الگوریتم}، \زیرخط{زمان اجرا} و \زیرخط{دلیل انتخاب خود} را نیز در ۴ الی ۵ خط به طور واضح توضیح دهید. همچنین در صورتی که زمان اجرای الگوریتم شما \زیرخط{سرشکن و یا تصادفی} است، آن را بیان کنید. در این سؤال نیاز به اثبات درستی الگوریتم و اثبات زمان اجرای آن نیست.

%ورودی برخی سؤال‌ها یک گراف است، اما برای حل مسئله ممکن است لازم باشد گراف جدیدی بسازید. در این موارد روش ساخت گراف را توضیح دهید.

% صورت اصلی مسئله را در این‌جا بنویسید.
\شروع{زیرمسائل}

\زیرمسئله[10]
شما در یک کتاب فروشی کار می‌کنید. جدیداً مشتری‌ها به جای این‌که کتاب را با مشخصات آن پیدا کنند از شما می‌پرسند که چند تا کتاب دارید که در بازۀ قیمتی $[a,b]$ هستند و پنج تا از این کتاب‌هایی که در این بازۀ قیمتی هستند را به آن‌ها معرفی کنید. یک داده‌ساختار مناسب برای نگه‌داری مشخصات کتاب پیشنهاد دهید که اضافه و کم کردن کتاب و همچنین پاسخ به پرسش‌های از نوع بالا را در سریع‌ترین زمان انجام دهد.

\پاسخ{}


\زیرمسئله[10]
گراف $G$ یک گراف جهت‌دار و همبند قوی است که دارای $n$ رأس و $2n+1$ یال است. الگوریتمی ارائه دهید که در صورت امکان یک یال از $G$ حذف کند که همچنان همبند قوی باقی بماند.

\پاسخ{
همیشه چنین یالی وجود دارد. از رأس دلخواه $v$ یک‌بار در جهت یال‌ها و یک بار در جهت‌عکس یال‌ها یکی از الگوریتم‌های DFS یا BFS را برای پیمایش گراف اجرا می‌کنیم. بنابر همبند قوی بودن گراف در هر دو حالت به تمام رأس‌ها می‌رسیم. از طرفی هر کدام از این اجراها یک درخت با $n-1$ یال تولید می‌کند که با روی هم قرار دادن یال‌های این دو درخت به حداکثر $2n-2$ یال می‌رسیم که اگر بقیه یال‌ها را حذف کنیم باز هم گراف همبند قوی می‌ماند.

}

\sloppy
\زیرمسئله[10]
یک لیست از رشته‌ها شامل یک «سه‌گانۀ ادغامی» $(a,b,c)$ است اگر $a$ و $b$ بتوانند با هم ادغام شوند، حروف آن‌ها در صورت لزوم جابه‌جا شود و رشتۀ $c$ حاصل شود. به‌طور مثال 
$(\text{''راستی``}, \text{''یاور``}, \text{''ویراستاری``})$ یک سه‌گانۀ ادغامی است ولی 
$(\text{''سالم``}, \text{''تقلا``}, \text{''استقلال``})$
چنین نیست. در مجموع $n$ رشته داریم که طول هر کدام نیز می‌تواند بین $1$ تا $n$ باشد. هر رشته، از حروف فارسی (بدون اعراب) تشکیل شده است (فاصله‌ها را حرف درنظر نمی‌گیریم و می‌توانیم فاصله‌ها را آزادانه اضافه و یا کم کنیم). یک الگوریتم کارا برای پیدا کردن تمام سه‌گانه‌های ادغامی رشته‌های داده شده طراحی کنید.

%\زیرمسئله[?]
%گراف همبند و بدون جهت $G$ را داریم که $n$ رأس و $n$ یال دارد. الگوریتمی ارائه دهید که یال‌های گراف را طوری جهت‌دهی کند که درجه خروجی هر رأس دقیقاً یک باشد.
%
%\پاسخ{
%	
%}

\زیرمسئله[10]
گراف وزن‌دار $G$ و رأس $s$ از آن داده شده‌اند. وزن یال‌ها در این گراف می‌تواند منفی یا مثبت باشد. این گراف دقیقاً یک دور مثبت دارد (ولی ممکن است چندین دور منفی داشته باشد). الگوریتمی با زمان $O(n^3)$ پیشنهاد دهید که تمام رئوسی که طولانی‌ترین مسیر از آن‌ها کمتر از $n$ رأس دارد را خروجی دهد.

\پاسخ{
%All simple longest paths will traverse fewer than |V| edges, whereas longest paths which are non-simple will contain a positive weight cycle. Negate edge weights in G and run Bellman-Ford from s to determine whether the positive weight cycle is reachable from s. If not, all vertices reachable from s satisfy the condition as all longest paths traverse fewer than |V | edges. Otherwise, Bellman-Ford finds all vertices reachable from s as well as finding the positive weight cycle reachable from s. Any vertex reachable from arbitrary vertex v on the cycle will not have a longest path from s traversing fewer than |V | edges, so run BFS or DFS to find the set of vertices reachable from v. Then, return the set of vertices reachable from s but not reachable from v. Bellman-Ford runs in O(|V ||E|) = O(|V |3) time and traversal takes O(|V | + |E|) = O(|V |2) time, so this algorithm runs in O(|V |3) time.
%Rubric:
%• 9 points for a correct algorithm. Common approach:
%– 3 points for algorithm to find positive weight cycle – 3 points for finding vertices reachable from cycle – 3 points for returning correct set
%• 2 points for correct running time analysis • 4 points if algorithm is O(|V |3)
%• Partial credit may be awarded
%
	
}


\پایان{زیرمسائل}